% Fragmento integrable. Preambulo anfitrion: amsmath, amssymb, longtable, hyperref.
\section{Aislamiento del límite y escalado del selector global}
\label{hmtfg:fragmento:aislamiento-del-limite-y-escalado-del-selector-global}

\subsection{Continuación indefinida de la primera raíz y exclusión del intervalo intermedio}
\label{hmtfg:escalado:15}

La clasificación completa demostrada anteriormente conserva la familia y el selector de primera raíz nueva. El transporte ordenado conserva los ceros y su mínimo; la palabra de retornos conserva la lectura cronológica nonádica.

\subsubsection{Clasificación y existencia a toda profundidad}
\label{hmtfg:escalado:15.1}

Escribamos \(W_n=\tau_1\cdots\tau_{2^n-1}\), con \(\tau_j=(-1)^{v_2(j)}\), y \(a_*=\min\Lambda_\tau\). La existencia y compacidad de \(\Lambda_\tau\) están demostradas en el desarrollo de prolongación, \ref{hmtfg:prolongacion:6.1}–\ref{hmtfg:prolongacion:6.2}; su mínimo existe por compacidad y orden.

\textbf{Clasificación.} Todo mapa unimodal estricto de máximo, sobre un intervalo invariante, cuyo crítico tiene período exacto \(2^n\) y cuyo itinerario \(K\) satisface \(K<\tau\), tiene itinerario \(W_n0\).

La inducción se realiza mediante el retorno efectivo. Para período al menos ocho, la comparación con \(\tau\), junto con los intervalos que atraparían órbitas de signos incompatibles, fuerza el prefijo \((+,-,+,+,+,-)\). De aquí se obtienen
\[
J_1=[d_2,d_4],\quad F(J_1)=[d_3,1],\quad
d_4<d_3,\quad F^2(J_1)=J_1.
\]
El retorno normalizado por \(h(x)=d_2x\) tiene máximo uno, período mitad e itinerario \(K'_j=-K_{2j}<\tau\). El orden se conserva porque
\[
O_{2q-1}(K)=-O_{q-1}(K'),\qquad
K_{2q}-\tau_{2q}=-(K'_q-\tau_q).
\]
Los casos iniciales son \(+0\) y \(+-+0\); la inducción reconstruye \(W_n0\). Para los mapas pares, también \(0<d_4<-d_2<1\), de modo que el retorno se extiende a \([-1,1]\) y coincide con \eqref{hmtfg:escalado:eq:6}.

Antes de \(a_*\), todos los itinerarios son menores que \(\tau\): los conjuntos de comparación estricta son abiertos por la separación demostrada en \ref{hmtfg:prolongacion:6.1}; el intervalo anterior a \(a_*\) es conexo y comienza con \(+^\infty<\tau\).

\textbf{Existencia de cada primera raíz.} Supóngase construida \(\lambda_{n-1}<a_*\), y sea \(p=2^{n-1}\). Los ceros de \(c_p\) en \([\lambda_{n-1},a_*]\) forman un conjunto finito no vacío: la función es analítica y \(c_p(a_*)\ne0\). Sea \(z<a_*\) su último cero. En \((z,a_*]\), \(c_p\) tiene signo \(\tau_p\). Para \(g_\lambda=f_\lambda^p\),
\[
g_\lambda'(0)=0,\qquad
c_{2p}(\lambda)=c_p(\lambda)+O(c_p(\lambda)^2).
\]
Cerca de \(z\), por la derecha, \(c_{2p}\) tiene signo \(\tau_p\); en \(a_*\), tiene el contrario. Existe un cero de \(c_{2p}\) en \((z,a_*)\), y su período es exactamente \(2p\), pues allí \(c_p\ne0\). La finitud de los ceros proporciona la primera raíz nueva en \((\lambda_{n-1},a_*)\). El último cero auxiliar prueba existencia; la elección efectiva sigue siendo la primera raíz. La base es \(\lambda_1=1/u_0(1)<a_*\).

\subsubsection{Límite de la sucesión original}
\label{hmtfg:escalado:15.2}

\textbf{Teorema.} El selector original existe a todos los niveles y satisface
\begin{equation}
\boxed{K_{\lambda_n}=W_n0,\qquad
\lambda_n\nearrow a_*=\min\Lambda_\tau.}
\label{hmtfg:escalado:eq:50}
\end{equation}
\textbf{Demostración.} La existencia precedente da una sucesión creciente acotada por \(a_*\). La clasificación fija todos sus itinerarios. Sea \(b\le a_*\) su límite. Para cada \(p\), los términos suficientemente avanzados conservan los signos opuestos \(\tau_p,\tau_{2p}\). Taylor y la cota uniforme \(B_p=16^p(16^p-1)/15\) dan
\[
|c_{2p}(\lambda_n)-c_p(\lambda_n)|
\le\tfrac12 B_p|c_p(\lambda_n)|^2,\qquad
|c_p(\lambda_n)|>2/B_p.
\]
Al pasar al límite se conserva \(\operatorname{sign}c_p(b)=\tau_p\), con módulo al menos \(2/B_p>0\). Para todo \(p\), esto implica \(b\in\Lambda_\tau\); la minimalidad da \(b=a_*\).

El resultado resuelve existencia indefinida, conservación de la cascada simbólica y selección de su primer parámetro límite para la regla global original. La completación HMT transporta esta sucesión y su límite conservando orden y lectura.

\subsubsection{Certificado de exclusión del intervalo intermedio}
\label{hmtfg:escalado:15.3}

El programa de puente certifica
\begin{equation}
\Lambda_\tau\cap[\lambda_{8,\rm inf},1.874038]=\varnothing,
\label{hmtfg:escalado:eq:51}
\end{equation}
donde
\[
\lambda_{8,\rm inf}
=1.874027744154450672897007573595092408435133414523352839687903472138975
\]
es el extremo inferior racional del intervalo certificado de \(\lambda_8\).

La cobertura contiene 374 intervalos exactamente adyacentes y ninguna región pendiente: 319 con \(c_{512}>0\), 10 con \(c_{1024}<0\), dos con \(c_{2048}>0\), todos opuestos al signo de \(\tau\), y 43 próximos a centros de períodos 256, 512 o 1024, excluidos por su segundo retorno.

En este último caso, \(g=f_\lambda^p\), \(d=g(0)\) y una cota uniforme \(|g''|\le M\) entre cero y \(d\) satisfacen \(M|d|<2\). Entonces
\[
|g(d)-d|\le\tfrac12M|d|^2<|d|\quad(d\ne0).
\]
Ambos retornos tienen el mismo signo; para \(d=0\), ambos son cero. Las dos situaciones excluyen \(\tau_{2p}=-\tau_p\) con signos estrictos.

La evaluación utiliza momentos racionales, colas geométricas del potencial y de su segunda derivada, y \(|u_0'''|<6\). El centrado paramétrico intersecta la envolvente natural con
\[
c_j(\lambda_{\rm medio})+
\partial_\lambda c_j(I)\,(I-\lambda_{\rm medio});
\]
la derivada se propaga sobre la caja completa. La autoaplicación de \([-1,1]\) se demuestra antes de restringir a ella las envolventes.

El certificado completo conserva la cobertura y comprueba su teselado exacto, junto con cada desigualdad de signo o de Taylor.

\subsubsection{Localización conjunta}
\label{hmtfg:escalado:15.4}

El cruce local \(\lambda_*\) realiza \(\tau\) por sus retornos reales a toda profundidad, de modo que \(a_*\le\lambda_*\). Por \eqref{hmtfg:escalado:eq:50}, \(\lambda_8<a_*\); junto con \eqref{hmtfg:escalado:eq:51},
\begin{equation}
\boxed{1.874038<a_*\le\lambda_*<1.874039.}
\label{hmtfg:escalado:eq:52}
\end{equation}
La localización conserva el itinerario canónico y sitúa su primer parámetro límite en el mismo intervalo que el cruce estable.

\subsection{Aislamiento por primera salida y coincidencia de los límites}
\label{hmtfg:escalado:16}

Esta sección completa el aislamiento identificado en sección~\ref{hmtfg:escalado:15.4}. Conserva la familia HMT \eqref{hmtfg:escalado:eq:1}, la selección mínima \eqref{hmtfg:escalado:eq:50}, la carta \eqref{hmtfg:escalado:eq:5} y la hoja estable de secciones~\ref{hmtfg:escalado:5}--\ref{hmtfg:escalado:7}. La prueba utiliza el orden \(a_*\le\lambda_*\): basta excluir una separación hacia el lado negativo del cono expansivo.

\subsubsection{Cotas superior e inferior del transporte transversal}
\label{hmtfg:escalado:16.1}

Además de \eqref{hmtfg:escalado:eq:9}, el lema 2.1 de \href{https://arxiv.org/pdf/1410.3277}{Hertling–Spandl, sección 2} proporciona \(\|D\Phi\|<0.9\) en la misma bola de radio \(0.01\), donde
\[
\Phi_u(u,\nu)=u-\frac{U(u,\nu)-u}{3.669}.
\]
Por tanto,
\begin{equation}
\left|1-\frac{U_u-1}{3.669}\right|<0.9,\qquad
U_u<1+1.9(3.669)=7.9711.
\label{hmtfg:escalado:eq:53}
\end{equation}
El número \(3.669\) pertenece al precondicionador publicado de este certificado analítico; la familia generada \eqref{hmtfg:escalado:eq:1} conserva sus coeficientes previos.

Tomemos \(\kappa=0.21\). Para dos puntos de la bola unidos por una secante con
\(\Delta u<0\), \(\|\Delta\nu\|_1\le\kappa|\Delta u|\), la integración de los bloques de \(DR\) sobre el segmento da
\begin{equation}
4.4034|\Delta u|
\le|\Delta u'|\le8.0887|\Delta u|,\qquad \Delta u'<0,
\label{hmtfg:escalado:eq:54}
\end{equation}
\begin{equation}
\|\Delta\nu'\|_1
\le(0.756+0.719\kappa)|\Delta u|
=0.90699|\Delta u|
<\kappa(4.4034)|\Delta u|.
\label{hmtfg:escalado:eq:55}
\end{equation}
Aquí \(4.4034=4.521-0.560\kappa\) y \(8.0887=7.9711+0.560\kappa\). Así, el cono de secantes conserva orientación y se expande mientras el segmento permanezca en la bola.

\subsubsection{Reducción de una separación infinita a una región finita}
\label{hmtfg:escalado:16.2}

Supongamos \(a_*<\lambda_*\), y pongamos
\[
f_n=R^{4+n}f_{a_*},\qquad
s_n=R^{4+n}f_{\lambda_*},\qquad
(\Delta u_n,\Delta\nu_n)=f_n-s_n.
\]
La cota positiva de \(u'\) y \(\|\nu'\|_1/u'<0.144386<\kappa\) sobre \(\Gamma(J)\) sitúan la secante inicial en el cono negativo. La hoja estable está escrita respecto de \(g\), con
\begin{equation}
s_n=g+e_n,\quad
\|e_n\|_{\mathcal A}<0.001666(0.83996)^n,\quad
|(e_n)_u|\le0.16\|(e_n)_\nu\|_1.
\label{hmtfg:escalado:eq:56}
\end{equation}
Los controles de entrada implican
\begin{equation}
|\Delta u_0|
<0.004993128+0.00004
 +0.16(0.001396077+0.00004)
=0.00526290032< h,\qquad h=0.006.
\label{hmtfg:escalado:eq:57}
\end{equation}
Mientras \(|\Delta u_n|<h\), \eqref{hmtfg:escalado:eq:55}–\eqref{hmtfg:escalado:eq:56} dan
\begin{equation}
\|f_n-g_0\|_{\mathcal A}
<(1+\kappa)h+0.001666+0.00004
=0.008966<0.01.
\label{hmtfg:escalado:eq:58}
\end{equation}
El punto \(s_n\) y el segmento entre ambos también permanecen en esa bola convexa. La expansión inferior \eqref{hmtfg:escalado:eq:54} obliga a una primera salida finita \(N\), y la cota superior aplicada al paso precedente proporciona
\begin{equation}
0.006\le-\Delta u_N<0.0485322,\qquad
\|\Delta\nu_N\|_1\le0.21|\Delta u_N|.
\label{hmtfg:escalado:eq:59}
\end{equation}
Se utiliza la desigualdad \(8.0887h=0.0485322\); el salto completo, y no únicamente la superficie de primera salida, queda incluido.

En consecuencia, \(f_N\) pertenece a la siguiente región:
\begin{equation}
\begin{split}
\mathcal C_-=\{\,g+(d_u,d_\nu)+e:\;&
-0.0485322\le d_u\le-0.006,\\
&\|d_\nu\|_1\le0.21|d_u|,\quad
|e_u|+\|e_\nu\|_1\le0.001666,\\
&|e_u|\le0.16\|e_\nu\|_1\,\}.
\end{split}
\label{hmtfg:escalado:eq:60}
\end{equation}
Como \(a_*\) realiza \(\tau\), todos sus retornos normalizados son autoaplicaciones unimodales de \([-1,1]\) con el mismo itinerario. Esta propiedad procede de la inducción de retornos, independientemente de la distancia de \(f_N\) a \(g_0\). Por eso el certificado siguiente trata \(\mathcal C_-\) intersectada con esa clase de autoaplicaciones.

\subsubsection{Exclusión certificada de la región de primera salida}
\label{hmtfg:escalado:16.3}

El certificado racional de exclusión por primera salida, reproducido en el suplemento computacional, cubre exactamente el intervalo de \eqref{hmtfg:escalado:eq:59} mediante 64 bandas racionales adyacentes. Todas quedan excluidas; ninguna subdivisión adicional ni región pendiente intervienen en el resultado.

De las 64 bandas, 54 presentan un signo contrario y diez un retorno contractivo de período 16 o 32. La mayor cota de Lipschitz de esos retornos es menor que \(0.475146491\); el horizonte máximo de exclusión es el iterado 64.

La evaluación conserva como variables compartidas los primeros cuarenta coeficientes de la perturbación, junto con una cota explícita de su cola infinita. Para cada banda, el centro es \(g_0\) desplazado en la coordenada \(u\). Si \(a,b\ge0\) acotan las sensibilidades duales en \(u,\nu\), el soporte de la incertidumbre del punto fijo y del residuo estable es
\begin{equation}
0.00004\max(a,b)
+0.001666\max\left(b,\frac{0.16a+b}{1.16}\right).
\label{hmtfg:escalado:eq:61}
\end{equation}
La segunda expresión se obtiene maximizando \(a|e_u|+b\|e_\nu\|_1\) bajo las dos restricciones de \eqref{hmtfg:escalado:eq:56}. De este modo se utiliza la orientación de la hoja estable demostrada previamente, además de su norma.

Sea \(x_j\) la órbita del centro y \(L_j\) su sensibilidad lineal a todos los coeficientes compartidos. El modelo de Taylor tiene la forma
\[
c_j=x_j+L_j(\Delta v)+r_j,\qquad |r_j|\le E_j.
\]
Si \(\rho_j\) acota \(|L_j(\Delta v)|+E_j\), la recurrencia exterior empleada es
\begin{equation}
E_{j+1}\le |f_0'(x_j)|E_j
+\tfrac12\sup|f_0''|\rho_j^2
+\sup|\Delta f'|\rho_j.
\label{hmtfg:escalado:eq:62}
\end{equation}
Las derivadas se acotan sobre el segmento de los dos estados. La órbita central se verifica dentro de \([-1,1]\); las otras órbitas pertenecen a ese intervalo por la autoaplicación incluida en el dominio certificado. Los cálculos utilizan redondeo racional exterior, incluyendo las colas geométricas.

Cada banda se excluye por una de dos desigualdades:

\par\noindent\textbf{1.}\enspace Un valor crítico tiene signo estrictamente opuesto a \(\tau_j\).
\par\noindent\textbf{2.}\enspace Para \(p=2^k\), el retorno \(H=f^p\) tiene constante de Lipschitz \(L_p<1\) sobre el segmento entre \(0\) y \(H(0)\). Entonces
\[
|H(H(0))-H(0)|\le L_p|H(0)|.
\]
Si \(H(0)\ne0\), los dos valores tienen el mismo signo; si \(H(0)=0\), ambos se anulan. En ambos casos queda excluido el itinerario estricto \(\tau_{2p}=-\tau_p\).

La constante \(L_p\) se obtiene propagando el segmento inicial \([-r,r]\), \(r\ge|H(0)|\), junto a las cajas de la órbita crítica y multiplicando las cotas exteriores de las derivadas. Se acotan los mapas de la banda completa, manteniendo \eqref{hmtfg:escalado:eq:61}.

\subsubsection{Teorema de coincidencia del parámetro de acumulación}
\label{hmtfg:escalado:16.4}

\textbf{Teorema.} Para la familia uniforme HMT \eqref{hmtfg:escalado:eq:1}, el límite de la selección original de primeras raíces coincide con el cruce estable certificado:
\begin{equation}
\boxed{\lambda_n\nearrow a_*=\lambda_*,
\qquad 1.874038<\lambda_*<1.874039.}
\label{hmtfg:escalado:eq:63}
\end{equation}

\textbf{Demostración.} \eqref{hmtfg:escalado:eq:50}–\eqref{hmtfg:escalado:eq:52} dan \(\lambda_n\uparrow a_*\le\lambda_*\), con ambos parámetros en \(J\). Si la desigualdad fuera estricta, \eqref{hmtfg:escalado:eq:54}–\eqref{hmtfg:escalado:eq:59} producirían un retorno \(f_N\in\mathcal C_-\) que realiza \(\tau\). La cobertura completa de sección~\ref{hmtfg:escalado:16.3} excluye precisamente esa posibilidad. Luego \(a_*=\lambda_*\).

La prueba controla todas las profundidades mediante una primera salida finita y una cobertura uniforme de su región. La identidad \eqref{hmtfg:escalado:eq:63} enlaza el límite del selector global con el cruce transversal.


\subsection{Identificación de los centros originales y escalado global}
\label{hmtfg:escalado:17}

El último paso conserva la selección de \eqref{hmtfg:escalado:eq:50}. La continuación local y los parámetros originales se comparan por su período exacto y su clase híbrida; la coincidencia se obtiene mediante una unicidad en un entorno paramétrico fijo.

\subsubsection{Realización analítica de grado dos y transversalidad}
\label{hmtfg:escalado:17.1}

El punto fijo \(g\) de sección~\ref{hmtfg:escalado:3} es par, analítico, de crítico cuadrático y combinatoria estacionaria de duplicación. El teorema 2.1 de \href{https://annals.math.princeton.edu/wp-content/uploads/annals-v164-n3-p01.pdf}{de Faria–de Melo–Pinto}, aplicado a esa única combinatoria, da un conjunto límite de un solo elemento, con extensión holomorfa propia de grado dos entre discos topológicos anidados. Su atracción sobre el intervalo identifica ese elemento con \(g\), pues \(Rg=g\).

El lema 3.2 de la misma fuente extiende un número fijo de retornos a una vecindad analítica completa de \(g\), con imágenes de grado dos y módulo positivo. La norma \eqref{hmtfg:escalado:eq:5} controla la norma uniforme sobre una vecindad compleja suficientemente pequeña de \([-1,1]\), porque \(|(x^2-1)/2.5|<1\) allí. La convergencia \eqref{hmtfg:escalado:eq:3} entra en dicha vecindad; la dependencia analítica de los retornos proporciona, para cierto entero fijo \(d\), una familia
\begin{equation}
\mathcal F_\lambda=R^d f_\lambda
\label{hmtfg:escalado:eq:64}
\end{equation}
de mapas de grado dos, analítica en un entorno complejo de \(\lambda_*\). Los discos y el número de retornos se fijan antes de variar \(\lambda\) en ese entorno.

La tangente paramétrica pertenece al cono expansivo, mientras la hoja estable satisface la pendiente de sección~\ref{hmtfg:escalado:5}. Los retornos preservan esa separación. La realización de grado dos identifica localmente la hoja estable con la clase híbrida del mapa infinitamente renormalizable; por ello \eqref{hmtfg:escalado:eq:64} es transversal a esa clase.

El cambio de espacio analítico se justifica detalladamente en \ref{hmtfg:clasificacion:8.2}, dedicada al transporte de la transversalidad. Una dirección tangente a la clase híbrida produciría tangentes renormalizadas decrecientes, por la convergencia uniforme sobre un disco híbrido y la estimación de Cauchy. La dirección certificada satisface, en cambio, \(|\dot f(1)|=|\dot u|/10\) y crece bajo los retornos por el cono expansivo. Esta evaluación común de ambos representantes prueba la transversalidad sin suponer equivalencia de sus normas.

\subsubsection{Unicidad en una vecindad fija}
\label{hmtfg:escalado:17.2}

Para \(n>d\), la clasificación de sección~\ref{hmtfg:escalado:15} da
\begin{equation}
K(\mathcal F_{\lambda_n})=W_{n-d}0.
\label{hmtfg:escalado:eq:65}
\end{equation}
El enderezamiento de un mapa de grado dos conserva la órbita crítica; \eqref{hmtfg:escalado:eq:65} identifica su clase híbrida con la del centro cuadrático canónico de período \(2^{n-d}\). Las cotas a priori de la sucesión real de duplicación son las empleadas en el teorema de universalidad de Lyubich.

El \href{https://www.maths.tcd.ie/EMIS/journals/Annals/149_2/lyubich.pdf\#page=69}{teorema 7.4 de Lyubich, pp. 387–388} da una sola intersección con cada una de esas clases, a niveles suficientemente altos, sobre la transversal analítica y dentro de una vecindad fija del parámetro de acumulación.

Por \eqref{hmtfg:escalado:eq:63}, \(\lambda_n\) entra finalmente en esa vecindad. La continuación local de sección~\ref{hmtfg:escalado:8}, indexada por el mismo período, pertenece a la misma clase híbrida. La unicidad implica
\begin{equation}
\boxed{\lambda_n=\widehat\mu_n\quad(n\ge n_0),}
\label{hmtfg:escalado:eq:66}
\end{equation}
donde \(\widehat\mu_n\) designa el centro local de período físico \(2^n\). Es una reindexación de \(\mu_j\) mediante un desplazamiento fijo determinado por los retornos y el blanco de sección~\ref{hmtfg:escalado:8}. La selección original de \(\lambda_n\) permanece intacta.

El entero \(n_0\) es existencial en esta prueba. La continuación de todos los niveles y la palabra \(W_n0\) ya están demostradas en \eqref{hmtfg:escalado:eq:50}; \eqref{hmtfg:escalado:eq:66} aporta la identidad métrica necesaria para el límite.

\subsubsection{Teorema de escalado para las primeras raíces HMT}
\label{hmtfg:escalado:17.3}

\textbf{Teorema.} Para la familia uniforme \eqref{hmtfg:escalado:eq:1}, sean \(\lambda_n\) las sucesivas primeras raíces nuevas de período crítico exacto \(2^n\). Existen \(C>0\), \(0<\theta<1\) y \(M<\infty\) tales que, para todo \(n\) suficientemente grande,
\begin{equation}
\lambda_*-\lambda_n=C\delta_{\mathrm F}^{-n}(1+e_n),
\qquad |e_n|\le M\theta^n,
\label{hmtfg:escalado:eq:67}
\end{equation}
y
\begin{equation}
\boxed{\lim_{n\to\infty}
\frac{\lambda_{n-1}-\lambda_{n-2}}
{\lambda_n-\lambda_{n-1}}=\delta_{\mathrm F}.}
\label{hmtfg:escalado:eq:68}
\end{equation}

\textbf{Demostración.} \eqref{hmtfg:escalado:eq:66} transporta el escalado local \eqref{hmtfg:escalado:eq:24} a las raíces originales; el desplazamiento finito de índices se absorbe en \(C,M\). Su signo es positivo porque \(\lambda_n<\lambda_*\).

El resto de potencia admite además una justificación directa mediante la holonomía de clases híbridas. El lema 7.3 de Lyubich da regularidad \(C^{1+\beta}\), con derivada no nula, sobre la transversal en el punto de acumulación. La dinámica unidimensional inestable es analítica y tiene multiplicador \(\delta_{\mathrm F}>1\); su coordenada linealizante hace que los centros sucesivos estén a distancia proporcional a \(\delta_{\mathrm F}^{-n}\). Componer con la inversa de la holonomía da un término lineal no nulo y un resto \(O(\delta_{\mathrm F}^{-(1+\beta)n})\). Así se puede tomar \(\theta=\delta_{\mathrm F}^{-\beta}\).

Escribiendo
\[
r_n=\frac{\delta_{\mathrm F} e_{n-1}-e_n}{\delta_{\mathrm F}-1},
\qquad K=\frac{M(\delta_{\mathrm F}+\theta)}{\delta_{\mathrm F}-1},
\]
se obtiene
\[
\lambda_n-\lambda_{n-1}
=C(\delta_{\mathrm F}-1)\delta_{\mathrm F}^{-n}(1+r_n),\qquad
|r_n|\le K\theta^{n-1}.
\]
Por tanto,
\begin{equation}
\left|
\frac{\lambda_{n-1}-\lambda_{n-2}}
{\lambda_n-\lambda_{n-1}}-\delta_{\mathrm F}
\right|
\le
\frac{\delta_{\mathrm F} K(1+\theta)\theta^{n-2}}
{1-K\theta^{n-1}},
\quad K\theta^{n-1}<1.
\label{hmtfg:escalado:eq:69}
\end{equation}
El lado derecho tiende a cero, demostrando \eqref{hmtfg:escalado:eq:68}.

\subsubsection{Encadenamiento y alcance del cierre}
\label{hmtfg:escalado:17.4}

La composición efectiva es: lector dodecafásico APP–TRIT–TPK; perfil \eqref{hmtfg:escalado:eq:1}; transporte ordenado de retornos \eqref{hmtfg:escalado:eq:38}–\eqref{hmtfg:escalado:eq:40}; clasificación y existencia de las primeras raíces \eqref{hmtfg:escalado:eq:50}; localización \eqref{hmtfg:escalado:eq:52}; aislamiento por primera salida \eqref{hmtfg:escalado:eq:63}; identidad eventual de centros \eqref{hmtfg:escalado:eq:66}; escalado \eqref{hmtfg:escalado:eq:67}–\eqref{hmtfg:escalado:eq:69}. La representación de Jacobi \eqref{hmtfg:escalado:eq:43} conserva el perfil completo de este lector y sus marcas permanecen en el antecedente.

La expansión espacial del lector HMT y el refinamiento nonádico aportan la genealogía y la evaluación a profundidad arbitraria. El escalado entre parámetros se demuestra mediante la dinámica del retorno, la transversalidad y el aislamiento que acabamos de componer. Las pruebas analíticas citadas intervienen sobre la familia ya construida; sus valores de comparación no seleccionan sus fases ni sus coeficientes.

El cierre se refiere al selector global original del sector uniforme \(\eta=0\), con las dependencias teoremáticas enumeradas. La ley asintótica conserva sus cuantificadores existenciales. La anchura de una caja de raíces finitas mide precisión de evaluación; la estimación \eqref{hmtfg:escalado:eq:69} expresa la convergencia hacia el límite.


Los programas y certificados del suplemento de reproducción permiten reproducir las partes finitas de la prueba. La continuación a toda profundidad se sustenta en las inducciones, el argumento uniforme de primera salida y los teoremas de identificación, con sus dominios explícitos.
